\section{Related work}\label{sec:related}

\paragraph{Interpretable and post-hoc XAI for CPS intrusion detection}
Some detectors expose interpretable internals: graph-attention weights that point to a deviating sensor \citep{deng2021gdn,tang2023grn}, attention on ICS traffic \citep{eltomy2024icps}, rules distilled from a logical neural network \citep{xu2024alrl}, or domain-aware relevance propagation \citep{rasheed2026xaicps}. Most of the literature applies SHAP or LIME to a detector or classifier, for SCADA and gas-pipeline autoencoders \citep{huong2022fedex,ha2022ifac}, federated and correlation-based IIoT detectors \citep{yazdinejad2025ppfl,birihanu2024correlation,abudurexiti2025cctak}, microgrids \citep{trivedi2025microgrid}, process-oriented attacks \citep{kenmogne2024netsoft,kenmogne2025intellisec} and network-traffic IIoT detectors \citep{oseni2023its,khan2022iiot,shoukat2025trustmyids,alani2023e2i3ds,attique2024iiot,ahakonye2024iotm,saiyed2026eshield}, or uses surrogate and example-based explanations \citep{zolanvari2023trust,panesar2026mmdcritic}. Interpretability is mostly validated qualitatively or by expert agreement. Several studies evaluate time-series data with random proportional splits or $k$-fold cross-validation (Supplementary Table~S1), a known source of optimistic bias \citep{arp2022dos,kaufman2012leakage}.

\paragraph{Evaluating explanations of intrusion detectors}
\citet{warnecke2020evaluating} proposed security-specific criteria (accuracy, sparsity, stability, efficiency, robustness, completeness), and \citet{wei2023xnids} and \citet{arreche2024exai} evaluate network-IDS explanations along similar axes. \citet{oinonen2026substations} compare seven XAI methods for a substation IDS, and WaXAI \citep{mathuros2024waxai} measures overlap between top-ranked features and attacked SWaT devices. \citet{fung2024attributions} provide the most systematic ground-truth evaluation: across 156 manipulations on SWaT, WADI and a Tennessee Eastman simulation, attributions at practical detection timings are often inaccurate, and an ensemble of raw-error, saliency and LEMNA \citep{guo2018lemna} attributions helps. We build on their protocol, change the unit of attribution to the control loop, follow explanations over the incident, and add timeliness.

\paragraph{Grouped, temporal and counterfactual explanations}
Owen values \citep{owen1977unions} extend Shapley values \citep{shapley1953value} to a priori unions of players, and sampling estimators make them tractable \citep{castro2009sampling,strumbelj2014explaining,lundberg2017shap}. ShaTS \citep{franco2026shats} groups time-series features by time step, sensor or SWaT process stage and reports execution time against a one-second budget; our loop game is an instance of such a priori grouping, with control loops as groups and an incident-local reference. \citet{mishra2026conditional} compute sensor--time attributions with conditional counterfactuals for root-cause analysis. AR-Pro \citep{ji2024arpro} produces diffusion-based counterfactual repairs with formal properties, evaluated on SWaT, WADI and HAI, and JoCE \citep{jeong2026joce} searches for sparse joint time--feature edits. Counterfactuals answer what would have to change rather than which loop is responsible; none of these methods has been evaluated against documented targets under deadlines.

\paragraph{Timeliness in CPS security}
\citet{cardenas2011attacks} and \citet{urbina2016limiting} analyse attacks and detection in terms of physical impact, and \citet{puccetti2026latency} argue that detection latency must be evaluated together with false-positive rates. GiBy \citep{venugopalan2025giby} explains detections through per-sensor safety bounds within milliseconds. To our knowledge, no prior work evaluates whether a \emph{correct} explanation reaches the operator within a process-derived deadline. \Cref{tab:positioning} summarises the closest work.

\begin{table}[t]
\centering\scriptsize
\caption{Positioning against the closest prior work (``--'': not reported).}\label{tab:positioning}
\resizebox{\linewidth}{!}{%
\begin{tabular}{@{}p{2.6cm}p{2.1cm}p{2.2cm}p{2.1cm}p{2.2cm}p{2.6cm}@{}}
\toprule
 & \citet{fung2024attributions} & ShaTS \citep{franco2026shats} & \citet{oinonen2026substations} & XAI-CPS \citep{rasheed2026xaicps} & \textbf{This work}\\
\midrule
Attribution unit & tag & tag / process stage & pixel & tag & control loop (+tag refinement)\\
Ground-truth localization & yes & case-based & no & expert agreement & yes (2 benchmarks, cross-checked)\\
At operational timing & yes (fails) & -- & n/a & -- & primary setting\\
Across alarm windows & no & no & no & no & yes (\TEA)\\
Physical prior in explainer & no & no & no & physics-informed LRP & yes (PI-CAX)\\
Physical deadline & no & no & no & no & yes\\
Statistical tests & ANOVA/\allowbreak Pearson & -- & -- & mean$\pm$sd & Friedman, Wilcoxon--Holm, bootstrap\\
\bottomrule
\end{tabular}}
\end{table}
